% !TEX program = lualatex
\documentclass[professionalfonts]{beamer}
\newcommand{\slidemode}{1}  % カラー投影モード
\usepackage{beamer_template}
\usepackage{enumerate}

\newcommand{\LectureNum}{8}
\newcommand{\LectureTitle}{周期 $2\pi$ の周期関数のフーリエ級数}
\newcommand{\CourseName}{応用数学}
\renewcommand{\TermName}{前期}

\begin{document}

\begin{frame}{第8回　周期 $2\pi$ の周期関数のフーリエ級数}
  \begin{block}{今回の内容}
  周期 $2\pi$ の周期関数のフーリエ級数
  \end{block}
  \vfill\hfill{\scriptsize \textcolor{gray}{第3章 フーリエ解析　§1}}
\end{frame}

\begin{frame}{定義：フーリエ展開（周期 $2\pi$）}
  \begin{block}{}
  $f(x)$ が周期 $2\pi$ の区分的に滑らかな関数であるとき、\\[2pt]
  次の無限級数に展開できる（フーリエ展開という）\\[2pt]
  \[
    f(x) = c_{0} + \sum_{n=1}^\infty (a_{n}\cos(nx) + b_{n}\sin(nx))
  \]
  \end{block}
\end{frame}

\begin{frame}{フーリエ係数の公式}
  \begin{block}{}
  \begin{align*}
    c_{0} &= \frac{1}{2\pi} \int_{-\pi}^{\pi} f(x)\,dx \\[6pt]
    a_{n} &= \frac{1}{\pi} \int_{-\pi}^{\pi} f(x)\cos(nx)\,dx \quad (n = 1, 2, 3, \ldots) \\[6pt]
    b_{n} &= \frac{1}{\pi} \int_{-\pi}^{\pi} f(x)\sin(nx)\,dx \quad (n = 1, 2, 3, \ldots)
  \end{align*}
  \end{block}
  \vfill
  {\small 計算でよく使う公式：$\sin(n\pi) = 0$，\;$\cos(n\pi) = (-1)^{n}$\quad（$n$：整数）}
\end{frame}

\begin{frame}{定理：三角関数の直交性}
  \begin{block}{}
  \begin{align*}
    \int_{-\pi}^{\pi} \cos(mx)\cos(nx)\,dx &= \begin{cases} 0 & (m \neq n) \\ \pi & (m = n) \end{cases} \\
    \int_{-\pi}^{\pi} \sin(mx)\sin(nx)\,dx &= \begin{cases} 0 & (m \neq n) \\ \pi & (m = n) \end{cases} \\
    \int_{-\pi}^{\pi} \cos(mx)\sin(nx)\,dx &= 0 \quad \text{（任意の $m, n$）}
  \end{align*}
  \end{block}
\end{frame}

\begin{frame}{問1}
  \begin{exampleblock}{問題}
    次の等式が成り立つことを証明せよ\\[2pt]
    $\displaystyle\int_{-\pi}^{\pi} \cos(mx)\cos(nx)\,dx = \begin{cases} 0 & (m \neq n) \\ \pi & (m = n) \end{cases}$
  \end{exampleblock}
\end{frame}

\begin{frame}{問1　解答}
  積和公式：$\cos(mx)\cos(nx) = \frac{1}{2}[\cos((m-n)x) + \cos((m+n)x)]$\\[4pt]
  $m \neq n$ のとき：各項 $\int_{-\pi}^{\pi}\cos(kx)\,dx = \left[\frac{\sin(kx)}{k}\right]_{-\pi}^{\pi} = 0$\\[4pt]
  $m = n$ のとき：$\cos^{2}(nx) = \frac{1}{2}(1+\cos(2nx))$\\[4pt]
  $\int_{-\pi}^{\pi}\cos^{2}(nx)\,dx = \frac{1}{2}\cdot 2\pi = \pi$
\end{frame}

\begin{frame}{直交性の証明（$\sin\cdot\sin$）}
  積和公式：$\sin(mx)\sin(nx) = \frac{1}{2}[\cos((m-n)x) - \cos((m+n)x)]$\\[4pt]
  $m \neq n$ のとき：各項 $\int_{-\pi}^{\pi}\cos(kx)\,dx = \left[\frac{\sin(kx)}{k}\right]_{-\pi}^{\pi} = 0$\\[4pt]
  $m = n$ のとき：$\sin^{2}(nx) = \frac{1}{2}(1-\cos(2nx))$\\[4pt]
  $\int_{-\pi}^{\pi}\sin^{2}(nx)\,dx = \frac{1}{2}\cdot 2\pi = \pi$
\end{frame}

\begin{frame}{直交性の証明（$\cos\cdot\sin$）}
  積和公式：$\cos(mx)\sin(nx) = \frac{1}{2}[\sin((m+n)x) + \sin((m-n)x)]$\\[4pt]
  任意の整数 $k$ に対して $\int_{-\pi}^{\pi}\sin(kx)\,dx = \left[-\frac{\cos(kx)}{k}\right]_{-\pi}^{\pi} = 0$\\[4pt]
  $m = n$ のときも：$\cos(nx)\sin(nx) = \frac{1}{2}\sin(2nx)$\\[4pt]
  $\int_{-\pi}^{\pi}\frac{1}{2}\sin(2nx)\,dx = 0$\\[6pt]
  よって任意の $m,\,n$ に対して $\int_{-\pi}^{\pi}\cos(mx)\sin(nx)\,dx = 0$
\end{frame}

\begin{frame}{定理：フーリエ展開の収束定理}
  \begin{block}{}
  \textbf{条件}：$f(x)$ は周期 $2\pi$ で 1 周期内において区分的に滑らか\\[6pt]
  \textbf{結論}：
  \begin{itemize}
    \item 連続点 $x_{0}$ では : フーリエ級数 $ \to f(x_{0})$
    \item 不連続点 $x_{0}$ では : フーリエ級数 $ \to (1/2)[f(x_{0}+0) + f(x_{0}-0)]$
  \end{itemize}
  \end{block}
\end{frame}

\begin{frame}{ギブス現象（p.\,79）}
  \begin{block}{}
  不連続点の近くでは、項数 $N$ を増やしても\\
  約 $9\,\%$ のオーバーシュートが残る（ギブス現象）
  \end{block}
  \vfill
  {\small 教科書 p.\,79 のグラフ（$N = 10, 20, 50$）参照}
\end{frame}

\begin{frame}{例題1　（p.\,78）}
  \begin{exampleblock}{問題}
    $f(x) = \begin{cases} 0 & (-\pi \leq x < 0) \\ x & (0 \leq x < \pi) \end{cases}, f(x+2\pi) = f(x)$ のフーリエ級数を求めよ\\[2pt]
  \end{exampleblock}
\end{frame}

\begin{frame}{例題1　解答（1/3）$c_{0}$ の計算}
  $-\pi \leq x < 0$ で $f(x) = 0$ なので積分区間を $[0,\pi)$ に縮小\\[6pt]
  \begin{align*}
    c_{0} &= \frac{1}{2\pi}\int_{-\pi}^{\pi} f(x)\,dx
     = \frac{1}{2\pi}\int_{0}^{\pi} x\,dx \\[4pt]
     &= \frac{1}{2\pi}\left[\frac{x^{2}}{2}\right]_{0}^{\pi}
     = \frac{1}{2\pi}\cdot\frac{\pi^{2}}{2}
     = \frac{\pi}{4}
  \end{align*}
\end{frame}

\begin{frame}{例題1　解答（2/3）$a_{n}$ の計算}
  \begin{align*}
    a_{n} &= \frac{1}{\pi}\int_{0}^{\pi} x\cos(nx)\,dx \\[4pt]
    &= \frac{1}{\pi}\left\{\left[\frac{x\sin(nx)}{n}\right]_{0}^{\pi} - \frac{1}{n}\int_{0}^{\pi}\sin(nx)\,dx\right\} \\[4pt]
    &= \frac{1}{\pi}\left\{\underbrace{\frac{\pi\sin(n\pi)}{n}}_{=\,0\;(\because\,\sin(n\pi)=0)} + \frac{1}{n}\left[\frac{\cos(nx)}{n}\right]_{0}^{\pi}\right\} \\[4pt]
    &= \frac{1}{n^{2}\pi}(\cos(n\pi) - 1)
     = \frac{(-1)^{n} - 1}{n^{2}\pi}
  \end{align*}
\end{frame}

\begin{frame}{例題1　解答（3/3）$b_{n}$ の計算}
  \begin{align*}
    b_{n} &= \frac{1}{\pi}\int_{0}^{\pi} x\sin(nx)\,dx \\[4pt]
    &= \frac{1}{\pi}\left\{\left[-\frac{x\cos(nx)}{n}\right]_{0}^{\pi} + \frac{1}{n}\int_{0}^{\pi}\cos(nx)\,dx\right\} \\[4pt]
    &= \frac{1}{\pi}\left\{-\frac{\pi\cos(n\pi)}{n} + \underbrace{\frac{1}{n}\left[\frac{\sin(nx)}{n}\right]_{0}^{\pi}}_{=\,0\;(\because\,\sin(n\pi)=0)}\right\}
     = \frac{-(-1)^{n}}{n}
     = \frac{(-1)^{n+1}}{n}
  \end{align*}
  \\[8pt]\textcolor{red}{\textbf{答}}：$f(x) = \dfrac{\pi}{4} + \displaystyle\sum_{n=1}^{\infty}\left[\frac{(-1)^{n}-1}{n^{2}\pi}\cos(nx) + \frac{(-1)^{n+1}}{n}\sin(nx)\right]$
\end{frame}

\begin{frame}{問2　（p.\,79）}
  \begin{exampleblock}{問題}
    $f(x) = x \;(-\pi \leq x < \pi),\; f(x+2\pi) = f(x)$ のフーリエ級数を求めよ\\[2pt]
  \end{exampleblock}
\end{frame}

\begin{frame}{問2　解答（1/2）$c_{0},\,a_{n},\,b_{n}$ の計算}
  $f(x) = x$ は奇関数なので $c_{0} = 0,\; a_{n} = 0$\\[6pt]
  \begin{align*}
    b_{n} &= \frac{1}{\pi}\int_{-\pi}^{\pi} x \sin(nx)\,dx \\[4pt]
    &= \frac{1}{\pi}\left\{\left[-\frac{x\cos(nx)}{n}\right]_{-\pi}^{\pi} + \frac{1}{n}\int_{-\pi}^{\pi}\cos(nx)\,dx\right\}
  \end{align*}
\end{frame}

\begin{frame}{問2　解答（2/2）$b_{n}$ の続き}
  $\displaystyle\int_{-\pi}^{\pi}\cos(nx)\,dx = \left[\frac{\sin(nx)}{n}\right]_{-\pi}^{\pi} = 0$\;（$\because \sin(n\pi) = 0$）より\\[6pt]
  \begin{align*}
    b_{n} &= \frac{1}{\pi}\left(-\frac{\pi\cos(n\pi)}{n} - \frac{(-\pi)\cos(-n\pi)}{n}\right) \\[4pt]
    &= \frac{1}{\pi}\cdot\left(-\frac{2\pi\cos(n\pi)}{n}\right)
     = \frac{-2(-1)^{n}}{n}
     = \frac{2(-1)^{n+1}}{n}
  \end{align*}
  \\[8pt]\textcolor{red}{\textbf{答}}：$f(x) = 2\displaystyle\sum_{n=1}^{\infty} \frac{(-1)^{n+1}}{n}\sin(nx) = 2\left(\sin x - \frac{\sin 2x}{2} + \frac{\sin 3x}{3} - \cdots\right)$
\end{frame}

\end{document}
